\documentclass[11pt]{article}
\usepackage{amssymb, amsmath, libertine, hyperref, xcolor, booktabs, enumitem, microtype}
\usepackage[T1]{fontenc}
\usepackage[margin=1in]{geometry}
\definecolor{dark-red}{rgb}{0.4,0.15,0.15}
\hypersetup{colorlinks, linkcolor=dark-red, citecolor=dark-red, urlcolor=dark-red}
\setlist[enumerate,1]{label=\textbf{\arabic*.}, leftmargin=*, itemsep=10pt}
\setlist[enumerate,2]{label=(\alph*), leftmargin=*, itemsep=4pt, topsep=4pt}
\setlist[enumerate,3]{label=(\roman*), leftmargin=*, itemsep=2pt, topsep=2pt}
\pagestyle{empty}

\newcommand{\prob}[1]{\medskip\noindent\textbf{\large #1}\par\smallskip}

\begin{document}

\begin{center}
{\Large\bfseries Math 111: Calculus \quad Homework 2}\\[4pt]
{Due Friday, October 9 by 10:00\textsc{p.m.} via Gradescope}
\end{center}

\medskip
\noindent\textbf{Covers} Active Calculus \S2.1--\S2.3, \S2.5, and \S2.6. Problem~1 can be
done now. Problems~2 and~4 use the chain rule (Monday, October~5), and Problem~3(c) uses
derivatives of inverse functions (Wednesday, October~7).

\smallskip
\noindent\textbf{Before you start.} Write in sentences and paragraphs, as an
explanation another student in this course could read and follow. Draw diagrams large
and label them. You are encouraged to work with classmates; cite everyone you worked
with, and write your solutions independently --- duplicated solutions are an Honor
Principle violation. Each problem is marked holistically on the 0--5 scale, with up
to two further points for the quality of the writing. If you need an extension, use
the \href{https://docs.google.com/forms/d/e/1FAIpQLSfIfgx44SFtUzefTNr-80yz35MVJaUqQttt96xVTU5mCQCOgg/viewform}{extension
form} \emph{before} the deadline.

\bigskip\hrule\bigskip

%%%%%%%%%%%%%%%%%%%%%%%%%%%%%%%%%%%%%%%%%%%%%%%%%%%%%%%%%%%%%%%%%%%%%%%%%%%%%%%%
\prob{1. A coffee cart, and where the money goes}

A coffee cart has been open for a few months. Let $N(t)$ be the number of cups it sells
per week, $P(t)$ the price of a cup in dollars, and $C(t)$ its total costs per week in
dollars, all $t$ weeks after opening. Ten weeks in, the owner knows:

\smallskip
\begin{center}
\begin{tabular}{@{}lll@{}}
\toprule
$N(10) = 800$ & $P(10) = 4.50$ & $C(10) = 2400$ \\
$N'(10) = -20$ & $P'(10) = 0.10$ & $C'(10) = 40$ \\
\bottomrule
\end{tabular}
\end{center}
\smallskip

\begin{enumerate}[resume]
\item[]
\begin{enumerate}
\item Let $R(t) = N(t)\,P(t)$ be the weekly revenue. Find $R(10)$ and $R'(10)$, with
units, and write one sentence saying what $R'(10)$ tells the owner.

\item The owner is surprised: the price is going \emph{up}, so why isn't revenue? Draw the
growing-rectangle picture from class for $R = N\cdot P$ at $t=10$, and use it to explain
the sign of $R'(10)$. Label which part of the picture is the effect of the price change
and which is the effect of the change in sales, and give the dollar amount of each.

\item How large would $P'(10)$ have to be for revenue to be holding steady at week~10, if
nothing else in the table changed? Then explain why ``if nothing else changed'' is
doing a lot of work in that question: which other number in the table would you expect
to move if the owner raised prices faster, and in which direction?

\item Let $A(t) = C(t)/N(t)$ be the cost of making one cup. Find $A(10)$ and $A'(10)$, with
units. Compare $A'(10)$ with $P'(10)$, and say what the comparison means for the money
the cart makes on each cup it sells.
\end{enumerate}
\end{enumerate}

%%%%%%%%%%%%%%%%%%%%%%%%%%%%%%%%%%%%%%%%%%%%%%%%%%%%%%%%%%%%%%%%%%%%%%%%%%%%%%%%
\prob{2. Climbing into the cold}

A hiker climbs a mountain. Let $a(t)$ be her altitude in feet, $t$ hours after she
starts, and let $T(a)$ be the air temperature in degrees Fahrenheit at altitude $a$
feet. Suppose that on this mountain, today, $T'(a) = -0.0035$ for every altitude $a$ on
her route. The temperature \emph{she experiences} at time $t$ is $H(t) = T\big(a(t)\big)$.

\begin{enumerate}[resume]
\item[]
\begin{enumerate}
\item Say in a sentence what $T'(a) = -0.0035$ means, with units. Then suppose that at
some moment she is climbing at $1200$ feet per hour. How fast is the temperature she
experiences changing at that moment? Explain why multiplying is the right thing to do,
using the units and a sentence about what happens over the next few minutes of her
climb --- not by quoting the chain rule.

\item Her altitude is recorded every hour:

\smallskip
\begin{center}
\begin{tabular}{@{}lrrrrr@{}}
\toprule
$t$ (hours) & 0 & 1 & 2 & 3 & 4\\
$a(t)$ (feet) & 3000 & 4100 & 5300 & 6200 & 6700\\
\bottomrule
\end{tabular}
\end{center}
\smallskip

Estimate $H'(2)$, with units. There is no formula for $a$ or for $H$.

\item From the table alone, over which part of the hike does $H'$ appear to be
increasing, and over which part decreasing? Say what each means for the hiker, in a
sentence about how cold it is getting.

\item A friend computes $H'(2)$ as $T'(2)\cdot a'(2)$. Explain what is wrong with writing
$T'(2)$, and why, on \emph{this} mountain today, the friend's number happens to come out
right anyway. Then describe a situation --- in words, about the weather --- in which the
friend's method would give the wrong answer.
\end{enumerate}
\end{enumerate}

%%%%%%%%%%%%%%%%%%%%%%%%%%%%%%%%%%%%%%%%%%%%%%%%%%%%%%%%%%%%%%%%%%%%%%%%%%%%%%%%
\prob{3. Where the rules come from}

\begin{enumerate}[resume]
\item[]
\begin{enumerate}
\item We proved the power rule in class for positive whole-number exponents by expanding
$(x+h)^n$. Use the quotient rule to show that it also holds for negative whole-number
exponents: if $n$ is a positive whole number, then
$\dfrac{d}{dx}\big[x^{-n}\big] = -n\,x^{-n-1}$.

\item Now the exponent $\tfrac12$. Assume that $s(x) = \sqrt{x}$ is differentiable for
$x > 0$. Differentiate both sides of $\big(s(x)\big)^2 = x$ with respect to $x$, and solve
for $s'(x)$. Check that your answer is what the power rule predicts for $x^{1/2}$.

\item Suppose $f$ is differentiable and has an inverse function $f^{-1}$, and that
$f(3) = 5$ and $f'(3) = 4$. Find $\big(f^{-1}\big)'(5)$. Then draw a picture that
explains your answer: the graph of $f$ near $(3,5)$ with its tangent line, the graph of
$f^{-1}$ near $(5,3)$ with its tangent line, and the line $y=x$. Say in a sentence why
reflecting a line across $y = x$ turns its slope into the reciprocal.

\item The rules have hypotheses. Answer each of the following; give an example if one
exists, and explain why not if it doesn't.
\begin{enumerate}
\item Two functions $f$ and $g$, \emph{neither} of which is differentiable at $x = 0$,
whose product $f g$ \emph{is} differentiable at $x = 0$.
\item Two functions $f$ and $g$, both differentiable at $x=0$, with $f'(0) = g'(0) = 0$
but $(fg)'(0) \neq 0$.
\end{enumerate}
\end{enumerate}
\end{enumerate}

%%%%%%%%%%%%%%%%%%%%%%%%%%%%%%%%%%%%%%%%%%%%%%%%%%%%%%%%%%%%%%%%%%%%%%%%%%%%%%%%
\prob{4. A rule that isn't, and a rule that is}

A student submits the following.

\smallskip
\begin{quote}
\textbf{Claim.} For any differentiable function $f$,
$\dfrac{d}{dx}\Big[\big(f(x)\big)^2\Big] = 2f(x)$.

\textbf{Argument.} Write $u = f(x)$. The power rule says that the derivative of $u^2$ is
$2u$. Substituting back, the derivative of $\big(f(x)\big)^2$ is $2f(x)$.
\end{quote}
\smallskip

\begin{enumerate}[resume]
\item[]
\begin{enumerate}
\item Identify the \emph{first} sentence of the argument that is wrong, and say precisely
what is wrong with it.

\item Give a specific function $f$ for which the claim is false, and show that it is false
by computing both sides.

\item State the correct rule for $\dfrac{d}{dx}\Big[\big(f(x)\big)^2\Big]$, explain where it
comes from, and check it on your example from part (b).

\item Your corrected rule gives a second route to the product rule. Check that
\[
f(x)\,g(x) = \tfrac14\Big[\big(f(x)+g(x)\big)^2 - \big(f(x)-g(x)\big)^2\Big],
\]
then differentiate the right-hand side using only your rule from part (c) and linearity.
Show that the result is $f(x)g'(x) + g(x)f'(x)$.
\end{enumerate}
\end{enumerate}

\end{document}
